\section{Scope and source boundaries}
\label{sec:scope}

The proof combines statements of different kinds.  We record their logical
roles explicitly.

\begin{enumerate}[label=\textup{(\arabic*)}]
\item Theorem~\ref{thm:simple-wall-point-column} is an exact identity inside
  the Gu--Yu--Yu completed formal comparison for \(r_+<r_-\).  It concerns
  only the ambient point coordinate.  It is not used to compose the global
  cobordism.
\item Theorem~\ref{thm:ordinary-flop-point-row} is an intrinsic point-section
  identity on one fixed Lee--Lin--Qu--Wang continuation domain.  It does not
  assert full descendant invariance or a Gamma/Fourier--Mukai square for all
  classes.
\item Corollary~\ref{cor:simple-wall-rank} requires its named one-wall
  sectorial realization hypothesis.  The formal Gu--Yu--Yu decomposition and
  the extremal Shen--Shoemaker asymptotics do not by themselves identify the
  intrinsic large-radius point section in a common receiver.
\item Propositions~\ref{prop:incomplete-gamma} and
  \ref{prop:punctual-corner} are exact analytic and algebraic countermodels.
  They are not asserted to arise from smooth projective quantum connections.
  They show why the local receivers cannot simply be glued by formal
  monodromy, pairing, or localized support.
\item Proposition~\ref{prop:support-collapse} uses the global orbit-cylinder
  marking.  The wall contributions vanish because their attaching nodes land
  in intermediate fixed loci on which this particular class restricts to
  zero.  No claim is made that arbitrary insertions enjoy the same collapse.
\item Proposition~\ref{prop:gamma-ratio-reduction} proves the
  coefficientwise Gamma-ratio and balance identities for individual
  fixed-graph contributions.  Proposition
  \ref{prop:clutching-tail-holonomicity} strengthens this to regular
  holonomicity and tempered growth for each complete clutching tail, including
  arbitrary bubble trees.  It does not determine the finite connection matrix
  between adjacent tails.
\item Hypothesis~\ref{hyp:marked-threshold} asks only that those finite
  connection matrices fix the cyclic point row.  Aleshkin--Liu prove a
  balanced contour theorem for finite-dimensional linear abelian GLSMs; they
  do not prove this marked compatibility for nonlinear virtual fixed graphs.
  The rational two-tail example preceding the hypothesis shows that separate
  holonomicity and a zero wall term do not imply it.
\item The global proof continues only the finite cyclic module generated by
  the point row and its \(z\)-covariant derivatives.  It does not assume that
  the entire equivariant quantum source has finite rank or admits a common
  continuation.  No formal Novikov variable is evaluated at a nonzero complex
  number.
\item Cai enters only Proposition~\ref{prop:cubic-endpoint}.  Equations
  (9.1)--(9.4) reconstruct the formal block from his displayed matrices, and
  (9.5)--(9.8) independently compute the point coefficient.  The birational
  conditional transport theorem is unchanged if this endpoint lemma is replaced by any
  other detected primary packet.
\end{enumerate}

Birational cobordism can have intermediate coarse quotients with finite
quotient singularities \cite{Wlodarczyk,AKMW}.  The global proof does not form
their quantum D-modules.  It uses only the smooth projective equivariant
completion, the smooth extreme quotients, and the complete fixed loci in the
gauged theory.  This is why resolving and composing the intermediate walls is
unnecessary.

Theorem~\ref{thm:birational-point-primary} is conditional on
Hypothesis~\ref{hyp:marked-threshold}.  It concerns the point row on a formal
primary packet in the marked cyclic continuation postulated by that
hypothesis.
It is not a claim that all Stokes matrices, Gamma lattices, or derived
categories are birational invariants.  The incomplete-Gamma and punctual
countermodels rule out those stronger inferences.  What survives is exactly
one Boolean row, obtained from one globally marked class.

The computational artifact has a similarly narrow role.  It recomputes the
cubic block and indicial polynomial, the period recurrence, and the
projective grading identity.  It does not verify virtual localization or
marked threshold compatibility.  In particular, it cannot turn the
conditional birational theorem into an unconditional one.
